\documentclass[12pt]{article}
\usepackage[utf8]{inputenc}
\usepackage[T1]{fontenc}
\usepackage{hyperref}
\title{Tiene Language Resources for K-12 Curriculum Development}
\author{Compiled by Claude Sonnet 5 (claude-sonnet-5) for Andr\'as Kornai}
\date{}
\begin{document}
\maketitle

\section{Introduction}
Tiene (ISO 639-3: tii; also Tiini, Kitiene, Tende) is a Bantu language of the Democratic Republic of the Congo. Very little was found: the searches turned up only general reference pages and a Joshua Project listing.

\section{Primary Educational Resources}
None found.

\section{Linguistic and Reference Materials}
General encyclopedia entry: \url{https://en.wikipedia.org/wiki/Tiene_language}

\section{Religious and Community Resources}
Joshua Project reports a New Testament (2007), no complete Bible, and no audio Bible or Jesus Film: \url{https://joshuaproject.net/index.php/languages/tii}. Its people-group profile: \url{https://joshuaproject.net/people_groups/15478}

\section{Digital and Audio Resources}
None found.

\section{Curriculum Development Context}
No dictionary, grammar, primer or school programme was found. A curriculum would start nearly from scratch, and further searching in French and in DRC sources would be needed.

\end{document}